\documentclass[12pt,a4paper]{article}
\usepackage[utf8]{inputenc}
\usepackage[spanish]{babel}
\usepackage{graphicx}
\usepackage{listings} % Para mostrar tu código C elegante
\usepackage{xcolor}
\usepackage{hyperref}

\title{Implementación de Control de Potencia PWM en Microcontroladores de 8 bits}
\author{Bernard Calixto}
\date{Febrero 2026}

\begin{document}

\maketitle

\section{Introducción}
Este proyecto documenta el diseño e implementación de un dimmer para luminarias LED utilizando el microcontrolador PIC12F675. Se prioriza la optimización de recursos y el uso de herramientas Open Source en entornos Linux (Fedora).

\section{Definición del Problema}
El presente proyecto nace de la necesidad de proporcionar una iluminación dinámica y estética para un elemento ornamental sacro: una efigie de plata con arco ornamental y cintas de seda fucsia y blanca. 

El desafío técnico se divide en tres vertientes:
\begin{itemize}
    \item \textbf{Efecto Visual:} Lograr una transición de luz suave (fading) que realce el brillo del metal y la textura de la seda sin parpadeos perceptibles.
    \item \textbf{Eficiencia Lumínica:} Optimizar la posición de los emisores para evitar la pérdida de intensidad por atenuación, descartando la iluminación posterior en favor de un enfoque frontal o lateral que resalte el relieve de la plata.
    \item \textbf{Restricción de Hardware:} Implementar el control en un microcontrolador de dimensiones reducidas (PIC12F675) para que el circuito sea discreto y transportable durante las procesiones.
\end{itemize}

\section{Cronograma de Desarrollo}
\begin{itemize}
    \item \textbf{Fase 1:} Validación de hardware (Blink LED).
    \item \textbf{Fase 2:} Implementación de ADC para control analógico.
    \item \textbf{Fase 3:} Generación de señal PWM y pruebas de carga.
\end{itemize}

\end{document}